% Reciprocal mathematical comparison, not proof of a historical interpretation.
\section{Existence policy in the Boger comparison}
\label{sec:p13-boger-annex}

\begin{annexitem}{annex:p13.boger.existence}{thm:p13-boger-existence-bridge}{Nonemptiness and opposition}
\begin{proof}
The definitions immediately make \(\omega=\neg\alpha\) and
\(\varepsilon=\neg\iota\), so each pair has exactly one true member
for every interpretation, including empty \(S\).

Assume \(S\ne\varnothing\). Then \(\alpha\) is equivalent to
\(S\subseteq P\). The condition \(\iota\) says there exists
\(u\in S\cap P\). Its negation \(\varepsilon\) says no such \(u\)
exists, equivalently every \(u\in S\) lies outside \(P\).
Also \(\omega=\neg(S\subseteq P)\), which is equivalent to existence
of \(u\in S\setminus P\). These are the four ordinary extensional
conditions under the stipulated nonemptiness.

If both \(\alpha\) and \(\varepsilon\) held, choose \(u\in S\).
The first would put \(u\) in \(P\); the second would exclude it.
Thus they cannot both hold. For \(S=\{u,v\}\), \(u\ne v\), and
\(P=\{u\}\), neither holds, so joint falsity is possible.

If \(\iota\) is false, choose \(u\in S\). Disjointness gives
\(u\notin P\), hence \(\omega\) is true. Therefore \(\iota,\omega\)
cannot both fail. The same two-member example makes them both true.

Now let \(S=\varnothing\). The explicit nonemptiness conjunct makes
\(\alpha\) false, and the empty intersection makes \(\iota\) false.
Their defined negations \(\omega,\varepsilon\) are true. In contrast,
ordinary inclusion \(\varnothing\subseteq P\) is true, and
\(\varnothing\setminus P=\varnothing\) has no member. This proves
the stated difference without identifying this comparison semantics
with the whole historical ontology.
\end{proof}
\end{annexitem}
